% Hanzo Blockchain Platform: Interplanetary-Scale Commerce Infrastructure
\documentclass[11pt,twocolumn]{article}
\usepackage{shared/hanzocover}
\usepackage{amsmath, amssymb}
\usepackage{graphicx}
\usepackage{booktabs}
\usepackage{hyperref}
\usepackage{enumitem}
\usepackage{xcolor}
\usepackage{float}
\hypersetup{colorlinks=true,linkcolor=black,citecolor=blue,urlcolor=blue}

\title{Hanzo Blockchain Platform: Interplanetary-Scale Commerce Infrastructure
with Decentralized Compute, Payments, and Governance}
\author{
    Zach Kelling \\
    \textit{CTO, Hanzo AI (Techstars '17)} \\
    \texttt{zach@hanzo.ai}
}
\date{2018}

\begin{document}
\hanzocoverpage

\twocolumn[
\begin{@twocolumnfalse}
\maketitle
\begin{abstract}
We present the Hanzo Blockchain Platform, a modular infrastructure designed to
support commerce at planetary and interplanetary scale. The platform provides a
complete set of interconnected services---decentralized DNS, computation
marketplaces, governance, messaging, notifications, payments, and
exchange---unified by a network token that pays for all operations. Built on an
enhanced Stellar Consensus Protocol (SCP~2.0) with a block-DAG architecture
optimized for solid-state storage, the system functions as an ``Apache web server
for blockchain'': a set of flexible, composable building blocks from which
developers construct decentralized applications without committing to a single
chain or execution environment. We describe the consensus model, storage layer,
developer framework, and service architecture, and outline a path toward
Byzantine-fault-tolerant, latency-compensated operation across satellite networks
and beyond.
\end{abstract}
\end{@twocolumnfalse}
]

\section{Introduction}
\label{sec:intro}

The question motivating this work is simple: what does commerce look like in one
hundred years? At planetary scale, billions of simultaneous participants require
infrastructure that no centralized system can provide. At interplanetary scale,
multi-minute light-speed delays between Earth, Moon, and Mars colonies demand
consensus protocols that tolerate massive network partitions and latency
asymmetry.

Hanzo addresses these requirements with a layered blockchain platform whose
design principles are:

\begin{itemize}[leftmargin=1.1em]
    \item \textbf{Own the interconnect layers.} A single network token pays for
    DNS resolution, transaction fees, computation, and service-layer operations.
    \item \textbf{Flexible building blocks.} Developers compose services the way
    they compose Unix pipes---each piece does one thing well.
    \item \textbf{Federated sovereignty.} Each sidechain maintains its own
    governance and hardware requirements; the mainnet is simply one federation
    among many.
\end{itemize}

\section{Consensus: Stellar Consensus Protocol 2.0}
\label{sec:consensus}

The original Stellar Consensus Protocol (SCP) introduced federated Byzantine
agreement with quorum slices, enabling open membership without sacrificing
safety. Hanzo extends SCP with three improvements:

\begin{enumerate}[leftmargin=1.1em]
    \item \textbf{Delegated Proof of Stake integration.} Validators stake tokens
    to participate, aligning economic incentives with consensus correctness. Stake
    weight influences quorum slice formation, preventing Sybil attacks without
    requiring a closed validator set.
    \item \textbf{Free-form path finding.} Cross-chain asset transfers use a
    generalized path-finding algorithm that discovers routes across federated
    chains, enabling atomic swaps without pre-configured corridors.
    \item \textbf{Latency compensation.} For high-latency links (e.g., Earth--Mars
    at 4--24 minutes one-way), the protocol supports asynchronous ballot phases
    with extended timeout windows, preserving liveness under partition.
\end{enumerate}

\section{Storage Architecture}
\label{sec:storage}

The platform employs a block-DAG rather than a linear chain, enabling concurrent
block production and higher throughput. The storage layer is designed for modern
hardware:

\begin{itemize}[leftmargin=1.1em]
    \item \textbf{LSM-tree key-value store.} Optimized for SSD access patterns
    with write-ahead logging and compaction.
    \item \textbf{Per-node graph database.} Tracks trust relationships, asset
    paths, and governance linkages.
    \item \textbf{Sharded storage.} Data is partitioned across validator groups,
    with erasure coding for redundancy.
\end{itemize}

\section{Proof of Reputation}
\label{sec:reputation}

Hanzo introduces a peer-to-peer reputation system in which trust endorsements
form a directed graph. Using a ``Kevin Bacon'' traversal, the protocol
back-traces trust paths to evaluate whether an endorsement carries weight
proportional to the endorser's aggregate trust score. Consensus-driven blocklists
remove bad actors and spammers without centralized moderation.

\section{Proof of Health}
\label{sec:health}

A novel consensus-level audit mechanism automatically verifies node integrity.
Blockchain analytics tools are built into the lowest protocol layer, providing
publicly accessible health markers, quality-of-service metrics, and diagnostic
support for token ecosystems.

\section{Service Layer}
\label{sec:services}

The platform exposes a rich service layer, each component operating as an
independent module composable via the decentralized message queue:

\subsection{Decentralized DNS}
Nodes self-name via consensus-based resolution. The naming system provides
internet-existence services, redundant failover, and clustering---replacing
centralized DNS with a Byzantine-fault-tolerant alternative.

\subsection{Computation Marketplace}
Decentralized cloud functions support a flexible programming model with arbitrary
hardware tags (GPU, TPM, quantum). Developers publish and consume computation
through a marketplace with fee-based pricing in the network token.

\subsection{Secure Computing Environment}
Smart contracts execute within Intel SGX trusted platform modules, providing
completely isolated execution contexts with hardware-attested key security. A
tiered trust model allows varying levels of hardware verification.

\subsection{Governance}
Each sidechain governs itself---choosing consensus parameters, hardware
requirements, and upgrade schedules. The mainnet coordinates cross-chain
governance through federated voting.

\subsection{Messaging and Notifications}
A decentralized message queue and event system provides the connective tissue for
cross-chain operations, enabling developers to build chain-agnostic applications.
Push notifications are native to the protocol.

\subsection{Payments and Exchange}
A built-in decentralized exchange provides cross-blockchain liquidity on the
plasma layer. Any natively minted token can serve as payment via the built-in
exchange, and security tokens are supported natively.

\section{Developer Platform}
\label{sec:developer}

Hanzo provides a developer-first experience:

\begin{itemize}[leftmargin=1.1em]
    \item \textbf{Flexible object model} with automatic API generation.
    \item \textbf{Multi-lingual SDKs} supporting Go, Node.js, Python, and EVM.
    \item \textbf{Sidechain primitives} in the developer model, analogous to
    threading primitives in Go---developers spawn application-specific sidechains
    as naturally as goroutines.
    \item \textbf{GraphQL and SQL API layers} for drop-in compatibility with
    existing software drivers.
    \item \textbf{Data streaming} via WebSocket APIs with atomic cache
    invalidation.
    \item \textbf{Cloud function and smart contract package manager} for
    distributing reusable libraries across the ecosystem.
\end{itemize}

\section{Scaling Architecture}
\label{sec:scaling}

\begin{table}[H]
\centering
\small
\begin{tabular}{ll}
\toprule
\textbf{Layer} & \textbf{Target Throughput} \\
\midrule
Mainnet consensus & 10,000 TPS \\
Sidechain (Plasma) & 1,000,000 TPS \\
Cross-chain atomic swaps & Multi-chain simultaneous \\
CDN / data acceleration & 1,000,000 req/s \\
\bottomrule
\end{tabular}
\caption{Target throughput by architectural layer.}
\label{tab:throughput}
\end{table}

The sidechain architecture provides one sidechain per application, with
cross-chain atomic swap support enabling simultaneous token existence on
multiple chains.

\section{Security}
\label{sec:security}

\begin{itemize}[leftmargin=1.1em]
    \item \textbf{Quantum resistance.} Post-quantum cryptographic primitives are
    integrated at the lowest layer, with lattice-based signatures available for
    high-security contexts.
    \item \textbf{Multisig wallets.} Built into the protocol rather than
    implemented as smart contracts, reducing attack surface.
    \item \textbf{Full encryption layer.} Configurable security contexts enable
    encrypted yet shared data, with the ability to purge data for compliance.
    \item \textbf{Byzantine fault tolerance.} The system tolerates massive
    network loss and partitioning---a requirement for satellite-linked and
    interplanetary deployments.
\end{itemize}

\section{Business Model}
\label{sec:business}

The network token funds all platform operations. Fee structures are transparent:
transactions carry small fees generally absorbed by businesses, while token
holders transact for free. Institutional-grade compliance, managed data pipes
with guaranteed bandwidth, and forced upgrade paths ensure enterprise readiness.

\section{Conclusion}
\label{sec:conclusion}

The Hanzo Blockchain Platform provides the infrastructure necessary for commerce
at scales that existing systems cannot support. By combining an enhanced
federated consensus protocol with a modular service architecture, developer-first
tooling, and latency-tolerant design, the platform offers a credible path toward
decentralized commerce infrastructure that operates from a single city to the
solar system. The ``Apache web server for blockchain'' philosophy ensures that
every building block is useful independently while composing into systems of
arbitrary sophistication.

\end{document}
